\section{Factorial Context Model}

The study distinguishes three conditions:

\begin{description}
\item[S: stationary.] The child performs the target cognitive and regulation rule using touchscreen interaction while seated or otherwise stationary.
\item[E: embodied non-AR.] The cognitive rule is preserved but response requires controlled whole-arm, device-orientation, or body movement without camera-based spatial augmentation.
\item[A: embodied AR.] The same rule is executed with movement plus camera-mediated spatial anchoring or environmental search.
\end{description}

Let the context variable be

\[
C_t \in \{S,E,A\}.
\]

The contrast \(E-S\) estimates the effect of adding embodied motor demand while avoiding AR-specific visual and spatial processing. The contrast \(A-E\) estimates the incremental contribution of augmented environmental coupling beyond embodiment. The contrast \(A-S\) captures the combined package and is therefore descriptive unless the component contrasts are also interpretable.

The cognitive rule, challenge definition, target probability, nominal difficulty, reward schedule, and recovery criterion should be matched across conditions as closely as the interaction modality permits.

\subsection{Mechanistic hypotheses}

\textbf{H1-E.} Embodied non-AR interaction will increase movement exposure relative to S while preserving cognitive-control performance within a preregistered non-inferiority margin.

\textbf{H2-AR.} AR will alter spatial/environmental engagement relative to E. Any claimed AR-specific advantage requires evidence in the \(A-E\) contrast rather than only \(A-S\).

\textbf{H3-R.} Recovery latency may differ by context, but faster recovery in A is not interpreted as improved regulation unless cognitive demand, task difficulty, and safety burden are sufficiently matched.

\textbf{H4-T.} In a later longitudinal experiment, practice of the same regulation rule across multiple contexts may improve independent transfer compared with repeated practice in a single context. This is a separate transfer hypothesis and should not be inferred from a one-session crossover.
